\section{Related Work}
\label{related-work}

\subsection{Simultaneous end-to-end speech translation}
\looseness=-1

While speech translation was initially performed using cascaded systems combining automatic speech recognition (ASR), machine translation (MT) and text-to-speech synthesis (TTS)~\cite{wahlster2000verbmobil, nakamura2006atr}, it recently evolved in fully end-to-end systems~\cite{jia19_translatotron,lee-etal-2022-direct,jia22-translatotron2,audiopalm} reducing error propagation and enabling transfer of non-linguistic information such as the speaker voice identity or prosody to the generated speech. At first trained with auxiliary text of phoneme translation tasks~\cite{jia22-translatotron2, streamspeech}, most recent works~\cite{seamless,hibiki,seed_liveinterpret_2, google_s2st} train directly on simultaneous S2TT and S2ST tasks so they can use the predicted text translation as a scaffolding for speech generation at inference time. Among direct ST training methods, those who achieve better speech naturalness are duplex audio systems that require to build a simultaneous ST dataset. They either rely on a synthetic data generation pipeline which includes a fine word-level text-to-translation alignment method~\cite{hibiki, google_s2st} or use a text LLM to split text into semantic chunks (a few words) that are individually translated thus providing chunk-level translation alignment~\cite{seed_liveinterpret_2} before collecting human-annotated interpretation data for finetuning purposes. \ours{} removes most of the complexity from synthetic data generation as it only requires sentence-level translation alignment easily obtained from punctuation. Thanks to an efficient RL process, it is then possible to reduce the translation latency of the model so it achieves state-of-the-art quality/latency trade-off in multiple input languages.


\subsection{Self-improvement of real-time translation systems}
\looseness=-1

RL methods to improve simultaneous translation systems were first explored in the context of text translation. Some works used preference-based approaches~\cite{simulpl} with preferences established in the context of simultaneous ST by prompting a text LLM while others applied online reinforcement procedures~\cite{simulpl, seqpo_simt} with sequence-level rewards as a combination of translation quality and latency metrics. Because they lack sub-sentence granularity in their preference or reward signals, it is difficult for these methods to find an appropriate balance between translation quality and latency during the RL process. More recently, Seed LiveInterpret 2.0~\cite{seed_liveinterpret_2} applied PPO~\cite{ppo} with a combination of intermediate evaluations of the generated sequences (\textit{process rewards}) and overall evaluation of the translation (\textit{outcome rewards}). Starting from a base supervised ST model trained with chunk-level alignment and finetuned on high-quality human interpretation data, they managed to strictly improve the quality/latency trade-off through RL. However, due to complex interactions between the numerous rewards they introduced, they encountered stability issues, reward hacking and had to rely on two different stages of RL training, using only outcome rewards at first before adding process rewards. On the other hand, \ours{} uses a single and straightforward reward system based on BLEU score~\cite{bleu} coupled with GRPO~\cite{grpo} without KL regularization as previously done by~\citet{magistral} to reduce memory requirements during training. Most importantly, it does not rely on any human interpretation or annotated data to finetune the model before reinforcement. On multilingual simultaneous ST tasks, \ours{} achieves state-of-the-art translation quality, latency, naturalness and speaker identity preservation. \ours{} is even able to adapt to a new input language after a light finetuning.